\documentclass[11pt]{article}

% -------------------------
% Packages
% -------------------------
\usepackage[a4paper,margin=1in]{geometry}
\usepackage{amsmath, amssymb, amsfonts}
\usepackage{physics}
\usepackage{bm}
\usepackage{graphicx}
\usepackage{float}
\usepackage{hyperref}
\usepackage{listings}
\usepackage{xcolor}
\usepackage{caption}
\usepackage{subcaption}

% -------------------------
% Code listing style
% -------------------------
\lstset{
    basicstyle=\ttfamily\small,
    keywordstyle=\color{blue},
    commentstyle=\color{gray},
    stringstyle=\color{teal},
    breaklines=true,
    frame=single,
    numbers=left,
    numberstyle=\tiny,
}

% -------------------------
% Title info
% -------------------------
\title{Learning Green's Functions as Operator Kernels for Linear Boundary Value Problems}
\author{Saurabh Bhardwaj \\ \small Independent Research}
\date{\today}

% -------------------------
\begin{document}
\maketitle

% -------------------------
\begin{abstract}
Green's functions provide an explicit representation of inverse linear operators
associated with differential equations and boundary conditions. While classical
numerical methods compute solutions for specific forcing terms, learning the
Green's function itself enables a compact and interpretable representation of
the underlying operator. In this work, we investigate the feasibility of learning
Green's functions from input-output data using constrained machine learning models.
Focusing on the one-dimensional Poisson equation with Dirichlet boundary conditions,
we frame Green's function learning as an operator learning problem and analyze the
role of physics-based structural constraints such as symmetry, boundary conditions,
and regularity. This project serves as a controlled study of kernel learning for
linear physical systems and as a baseline for future extensions to higher-dimensional
and parameterized operators.
\end{abstract}

% =========================================================
\section{Introduction and Motivation}

Linear differential operators appear ubiquitously in physics, governing phenomena
such as electrostatics, diffusion, elasticity, and quantum mechanics. For a given
linear operator and boundary conditions, the solution to a boundary value problem
is fully characterized by the corresponding Green's function, which represents the
kernel of the inverse operator.

Classical numerical methods solve for the system response to a specific forcing
function, whereas the Green's function encodes the response to \emph{all} admissible
forcings. Learning the Green's function, therefore, constitutes a stronger and more
interpretable objective than learning individual solutions.

Recent developments in machine learning have focused on learning solution maps for
parametric partial differential equations. In contrast, this work explores the
learning of the operator kernel itself, explicitly incorporating physical structure
into the learning process.

% =========================================================
\section{Mathematical Background}

Consider a linear boundary value problem
\begin{equation}
L u(x) = f(x), \quad x \in \Omega,
\end{equation}
where $L$ is a linear differential operator acting on functions defined on a domain
$\Omega$, subject to appropriate boundary conditions.

If the inverse operator $L^{-1}$ exists, the solution may be written in integral form:
\begin{equation}
u(x) = \int_\Omega G(x,x') f(x') \, dx',
\end{equation}
where $G(x,x')$ is the Green's function associated with $L$ and the imposed boundary
conditions.

The Green's function is thus the integral kernel of the inverse operator and fully
encodes the system's response characteristics.

% =========================================================
\section{Problem Formulation}

\subsection{Physical setting}

We focus on the one-dimensional Poisson equation with Dirichlet boundary conditions:
\begin{equation}
-\frac{d^2 u}{dx^2} = f(x), \quad x \in [0,1],
\end{equation}
subject to
\begin{equation}
u(0) = u(1) = 0.
\end{equation}

The forcing function $f(x)$ is assumed to lie in $L^2([0,1])$.

\subsection{Green's function representation}

The corresponding Green's function satisfies
\begin{equation}
-\frac{d^2}{dx^2} G(x,x') = \delta(x-x'),
\end{equation}
with boundary conditions
\begin{equation}
G(0,x') = G(1,x') = 0.
\end{equation}

The analytic Green's function is given by
\begin{equation}
G(x,x') =
\begin{cases}
x(1-x') & x \le x', \\
x'(1-x) & x > x'.
\end{cases}
\end{equation}

\subsection{Learning objective}

Rather than learning the solution map $f \mapsto u$, we aim to learn an approximation
to the Green's function itself:
\begin{equation}
\hat{G}_\theta : [0,1] \times [0,1] \rightarrow \mathbb{R},
\end{equation}
such that
\begin{equation}
u(x) \approx \int_0^1 \hat{G}_\theta(x,x') f(x') dx'.
\end{equation}

This corresponds to learning an explicit approximation of the inverse operator
$L^{-1}$.

% =========================================================
\section{Structural Constraints}

The Green's function satisfies several structural properties:

\paragraph{Symmetry}
Since the operator is self-adjoint,
\begin{equation}
G(x,x') = G(x',x).
\end{equation}

\paragraph{Boundary conditions}
\begin{equation}
G(0,x') = G(1,x') = 0.
\end{equation}

\paragraph{Regularity}
The Green's function is smooth for $x \neq x'$, with a discontinuity in its first
derivative at $x=x'$.

These properties are enforced during training via constraint terms in the loss
function.

% =========================================================
\section{Data Generation}

Training data consists of pairs $(f_i, u_i)$ generated as follows:
\begin{itemize}
\item Sample forcing functions $f_i(x)$ from a predefined function family
(e.g.\ Fourier series with random coefficients or Gaussian random fields).
\item Compute solutions $u_i(x)$ using analytic convolution with the true Green's
function or high-resolution numerical solvers.
\end{itemize}

The Green's function itself is not provided to the model during training.

% =========================================================
\section{Baseline Methods}

The learned Green's function is compared against the following baselines:
\begin{itemize}
\item Finite difference solvers based on matrix inversion.
\item Spectral methods using Fourier representations.
\item Kernel ridge regression treating the solution operator as a black-box kernel.
\end{itemize}

These baselines establish numerical accuracy and highlight the trade-offs between
classical solvers and learned operator representations.

% =========================================================
\section{Learning Methodology}

A parameterized model $\hat{G}_\theta(x,x')$ is trained by minimizing the loss
\begin{equation}
\mathcal{L}(\theta) =
\mathbb{E}_{f} \left[
\left\| u_{\text{true}} -
\int_0^1 \hat{G}_\theta(x,x') f(x') dx'
\right\|_{L^2}^2
\right]
+ \lambda \mathcal{R}_{\text{physics}},
\end{equation}
where $\mathcal{R}_{\text{physics}}$ enforces symmetry, boundary conditions, and
smoothness.

The focus is on interpretability and physical consistency rather than architectural
complexity.

% =========================================================
\section{Evaluation Criteria}

Model performance is evaluated using:
\begin{itemize}
\item $L^2$ error between predicted and true solutions.
\item Error between learned and analytic Green's functions.
\item Generalization to unseen forcing functions.
\item Stability under perturbed or noisy inputs.
\end{itemize}

% =========================================================
\section{Scope and Limitations}

This study is restricted to linear operators on fixed domains with prescribed
boundary conditions. No claims are made regarding computational efficiency relative
to classical numerical solvers.

The primary objective is to establish a clean and interpretable framework for
learning operator kernels in physics-informed settings.

% =========================================================
\section{Future Extensions}

Natural extensions include:
\begin{itemize}
\item Higher-dimensional domains.
\item Parameterized operators.
\item Time-dependent Green's functions.
\item Nontrivial geometries and boundary conditions.
\end{itemize}

These extensions are left for future work.

% -------------------------
\end{document}
